\begin{proposition}[What trivial normalization pullback forces]\hypertarget{tag-01AZ}{}\label{cdp:st:A.5.1}
\StatementRecord{cdp:st:A.5.1}{proved}{01AZ}{\SeriesRef{cdp:st:A.1.1}, \SeriesRef{cdp:st:A.4.1}}

Let $S$ be a connected reduced full fibre of a proper map from a complex manifold to a smooth curve. For every line bundle $A$ on $S$ with $\eta^*A\simeq\OO_{\widetilde S}$, the ambient group contains the canonical image of $H^0(\widetilde S,\eta^*(N^*\otimes A))$. A local base coordinate and a choice of trivialization of $\eta^*A$ identify this subspace with $H^0(\widetilde S,\OO_{\widetilde S})$. In particular it is nonzero. This applies to the subcase of Proposition~7.3 in which all $\eta^*(L_t|_S)$ are trivial. It is compatible with the vanishing of $H^0(S,Q\otimes A)$ through the exact maps of \ref{cdp:st:A.1.1}, as \ref{cdp:st:A.4.1} demonstrates.
\end{proposition}


\begin{proof}
\begin{enumerate}[label={$\langle1\rangle$\arabic*.},leftmargin=*]
\item 
\noindent\textbf{Step 1.} Tensoring the canonical injection in (A.3) by $A$ and applying the projection formula gives
$\eta_*\eta^*(N^*\otimes A)\hookrightarrow K\otimes A\hookrightarrow V\otimes A$.
Taking global sections preserves injectivity and defines the asserted canonical subspace. A local coordinate of the base trivializes $N^*$; together with a choice of trivialization of $\eta^*A$, it identifies the source with $H^0(\widetilde S,\OO_{\widetilde S})$. Every connected component of the compact normalization has the nonzero constant function $1$, proving nonvanishing. No normal-crossings assumption is needed for this inclusion.

\end{enumerate}
\end{proof}